\chapter{FlashAttention：Online Softmax 的严格推导与 IO 复杂度}

\begin{warningBox}[避坑指南与核心警示]
\textbf{阅读前须知}：FlashAttention（2022）可能是近几年影响最大的一个推理/训练优化——它\textbf{不改变任何数学结果}，却让注意力快了数倍、显存从 $O(N^2)$ 降到 $O(N)$，从而让长上下文成为可能。它的核心是一个非常漂亮的代数技巧：\textbf{Online Softmax}。本篇把它从头推导一遍。

\textbf{前置}：第 2 篇（Softmax 平移不变性）、第 4 篇（注意力公式）、第 8 篇（存储层级、算术强度、算子融合）。

\textbf{配套实验}：\textbf{\texttt{lab09\_flash\_attention/flash\_attention\_numpy.py}}——用 NumPy 实现 Online Softmax、分块 FlashAttention 与 FlashDecoding 的 split-K 合并，逐一验证与朴素注意力完全相等；并在 GPU 上实测朴素注意力与 PyTorch SDPA 的显存差距。
\textbf{真机版}：\textbf{\texttt{lab09\_flash\_attention/cuda\_softmax/}}——把本节和下一节的 Online Softmax 写成\textbf{真正的 CUDA C++ kernel}，在你自己显卡上量出 DRAM 带宽、L2 命中与占用率。不需要装 CUDA toolkit。
\end{warningBox}

\vspace{0.8em}\hrule\vspace{0.8em}

\section{[低速] 一、朴素注意力慢在哪里？}

对序列长度 $N$、头维度 $d$，朴素实现分三步，每一步都是一个独立的 kernel：

\begin{center}
\small
\begin{tabularx}{\textwidth}{l X X}
\toprule
\textbf{标准 Attention 步骤} & \textbf{数学算子公式} & \textbf{朴素显存读写量（HBM IO）} \\
\midrule
1. 计算打分矩阵 & $S = Q K^T / \sqrt{d}$ & 读 $Q, K$ ($2Nd$) $\longrightarrow$ 写 $S$ 到显存 ($N^2$) \\
2. 概率归一化 & $P = \mathrm{softmax}(S)$ & 读 $S$ ($N^2$) $\longrightarrow$ 写 $P$ 到显存 ($N^2$) \\
3. 加权求和输出 & $O = P V$ & 读 $P$ ($N^2$), $V$ ($Nd$) $\longrightarrow$ 写 $O$ 到显存 ($Nd$) \\
\midrule
\textbf{总访存量} & \multicolumn{2}{l}{\textbf{高达 $4Nd + 2N^2$ 字节（受 $O(N^2)$ 显存访问严重瓶颈束缚）}} \\
\bottomrule
\end{tabularx}
\end{center}

\begin{itemize}
  \item \textbf{显存}：必须存下 $N \times N$ 的 $S$ 和 $P$。$N = 32\text{K}$ 时单个头就是 $10^9$ 个元素，BF16 下 2GB——而一层有几十个头。
  \item \textbf{访存}：总共读写 $\Theta(Nd + N^2)$ 个元素。而 $d$ 通常只有 64\textasciitilde{}128，$N^2$ 项完全主导。Softmax 这一步每个元素只做几次运算，算术强度极低（第 8 篇），\textbf{GPU 大部分时间在搬运 $N^2$ 的中间矩阵}。
\end{itemize}

\textbf{自然的想法}：像第 8 篇说的那样做算子融合——把 Q、K、V 分块读进片上 SRAM，算完直接输出 O，$S$ 和 $P$ 永远不落到显存。
\textbf{障碍}：Softmax 需要\textbf{整行}的最大值和总和，才能归一化任何一个元素。只看到一个块时，你不知道这一行后面还有没有更大的值。

\vspace{0.8em}\hrule\vspace{0.8em}

\section{[推导] 二、Online Softmax：一遍扫描算出 Softmax}

\subsection{传统的"安全 Softmax"要扫三遍}

对一行 $x_1, \dots, x_N$（第 2 篇）：
\begin{enumerate}
  \item 第一遍求 $m = \max_i x_i$；
  \item 第二遍求 $\ell = \sum_i e^{x_i - m}$；
  \item 第三遍输出 $p_i = e^{x_i - m} / \ell$。
\end{enumerate}

\subsection{把前两遍合成一遍}

维护两个"运行中"的量：到目前为止看过的最大值 $m_j$ 和对应的和 $\ell_j$：

\[
m_j = \max(m_{j-1}, x_j), \qquad \ell_j = \ell_{j-1} \cdot e^{m_{j-1} - m_j} + e^{x_j - m_j}
\]

初始 $m_0 = -\infty$，$\ell_0 = 0$。

\textbf{定理}：对所有 $j$，$\ell_j = \sum_{i=1}^{j} e^{x_i - m_j}$。

\textbf{证明（归纳法）}：
\begin{itemize}
  \item $j = 1$：$m_1 = x_1$，$\ell_1 = 0 + e^{x_1 - x_1} = 1 = \sum_{i=1}^{1} e^{x_i - m_1}$ [OK]
  \item 假设 $\ell_{j-1} = \sum_{i=1}^{j-1} e^{x_i - m_{j-1}}$，则
\end{itemize}

\[
\ell_j = \left(\sum_{i=1}^{j-1} e^{x_i - m_{j-1}}\right) e^{m_{j-1} - m_j} + e^{x_j - m_j} = \sum_{i=1}^{j-1} e^{x_i - m_j} + e^{x_j - m_j} = \sum_{i=1}^{j} e^{x_i - m_j} \quad \blacksquare
\]

关键就是那个\textbf{修正因子} $e^{m_{j-1} - m_j}$：当最大值变大时，之前累加的和是以旧的最大值为基准的，乘上这个因子就把它们"换算"到新基准下。因为 $m_j \ge m_{j-1}$，这个因子 $\le 1$，\textbf{永远不会溢出}。

上面是逐个元素更新；换成逐块更新（一次处理一个块 $B$）形式完全一样：

\[
m_{\text{new}} = \max\left(m_{\text{old}}, \max_{i \in B} x_i\right), \qquad \ell_{\text{new}} = \ell_{\text{old}} \, e^{m_{\text{old}} - m_{\text{new}}} + \sum_{i \in B} e^{x_i - m_{\text{new}}}
\]

\vspace{0.8em}\hrule\vspace{0.8em}

\section{[加速] 三、把输出也"在线"算出来：FlashAttention 的核心}

注意力的一行输出是 $o = \sum_i p_i v_i = \frac{1}{\ell} \sum_i e^{x_i - m} v_i$（$x_i$ 是这一行第 $i$ 个分数 $q \cdot k_i / \sqrt{d}$）。
我们连最后的归一化都不急着做，只维护一个\textbf{未归一化的累加器}：

\[
\text{acc}_j = \text{acc}_{j-1} \cdot e^{m_{j-1} - m_j} + \sum_{i \in B_j} e^{x_i - m_j} \, v_i
\]

\textbf{定理}：处理完第 $j$ 块后，$\text{acc}_j = \sum_{i \in B_1 \cup \dots \cup B_j} e^{x_i - m_j} v_i$。证明与上一节完全相同（把标量 1 换成向量 $v_i$）。
于是全部块处理完后：

\[
o = \frac{\text{acc}_{\text{final}}}{\ell_{\text{final}}}
\]

\textbf{这和一次性算完整的 Softmax 再乘 V 在数学上严格相等}，不是近似。

\subsection{FlashAttention 算法（前向，单个头）}

\begin{lstlisting}[language=Python]
把 Q 按行切成若干块 Q_1..Q_Tr（每块 Br 行），K、V 切成 K_1..K_Tc（每块 Bc 行）
for 每个 Q 块 Q_i（不同的 Q 块互相独立 → 分给不同的 SM 并行）:
    把 Q_i 读进 SRAM；初始化 m = -∞, l = 0, acc = 0（都在 SRAM / 寄存器里）
    for 每个 K/V 块 j:
        把 K_j、V_j 读进 SRAM
        S_ij = Q_i K_jᵀ / √d                    # Br × Bc 的小矩阵，只存在片上
        (因果掩码：若 K_j 整块都在 Q_i 之后，直接跳过；对角块内部再逐元素掩码)
        m_new = max(m, rowmax(S_ij))
        P_ij  = exp(S_ij - m_new)
        l     = l · exp(m - m_new) + rowsum(P_ij)
        acc   = acc · exp(m - m_new) + P_ij V_j
        m     = m_new
    O_i = acc / l，写回显存；另外存下 logsumexp = m + log l（训练的反向传播要用）
\end{lstlisting}

（这是 FlashAttention-2 的循环顺序：外层循环 Q，内层循环 K/V。FlashAttention-1 的外层是 K/V。）

\subsection{它省下了什么？}

\begin{itemize}
  \item \textbf{显存}：$S$、$P$ 从不写入显存，额外空间只有每行的 $m$、$\ell$：$O(N)$。
  \item \textbf{访存}：Q、O 各读写一次；K、V 对每个 Q 块读一次，共 $N / B_r$ 次。设片上 SRAM 大小为 $M$（按元素计），块大小 $B_r \sim M / d$，论文证明 FlashAttention 的 HBM 访问量为 $\Theta\left(\frac{N^2 d^2}{M}\right)$，而朴素实现为 $\Theta(Nd + N^2)$。
\end{itemize}

  两者之比的量级是 $\frac{M}{d^2}$——\textbf{这是一个与 $N$ 无关的常数}。
  按配套实验 3 的简化模型（约 100 KB SRAM 要同时放下 Q、K、V、O 四个块），$d = 64$ 约省 6 倍，$d = 128$ 约省 1.5 倍。
\begin{noteBox}[核心导读与背景]
注：实验 3 用的是\textbf{元素个数}口径，它数出的是 $2B_r/d = M/(2d^2)$，和论文的 $M/d^2$ 差一个 2。这不是矛盾——实验 3 的朴素项数了 4 趟 $N^2$（写 $S$、读 $S$、写 $P$、读 $P$），而论文的 $\Theta(Nd + N^2)$ 只按元素数算一趟。\textbf{量级和"与 $N$ 无关"这个结论都不受影响}，但引用具体倍数时要说清用的是哪套账。
另外，小 $N$ 时这个常数还到不了（尾块占比大）：$d=64$ 时要到 $N \gtrsim 16\text{K}$ 才收敛到 6 倍，$N=1024$ 时只有 4.9 倍。真实 kernel 的块形状与 L2 缓存命中也会让具体数字不同，但"常数倍、与 $N$ 无关"这一点不变。
\end{noteBox}

\begin{itemize}
  \item \textbf{计算量}：FLOPs 与朴素完全相同（甚至略多一点点 exp 修正）。
  \item \textbf{随 $N$ 增长的真正收益是显存}：朴素实现需要 $O(N^2)$ 的额外显存，FlashAttention 只要 $O(N)$。再加上多个 kernel 融合成一个、不再产生 FP32 的 $N^2$ 中间张量，实测加速往往比访存模型预测的还大（配套实验 5 在 5060 Ti 上 $N = 4096$ 时约 16 倍）。
\end{itemize}

\begin{tipBox}[核心直觉与思考]
{}[思考] \textbf{一个反直觉的教训}：FlashAttention 没有减少任何计算，只是\textbf{减少了数据搬运}，就换来了数倍加速。这正是第 8 篇 Roofline 的核心观点——在访存受限的区域，优化的对象是字节而不是 FLOPs。它的反向传播更进一步：\textbf{宁可重新计算 $S$ 和 $P$，也不去读存下来的 $N^2$ 矩阵}，因为重算比读显存更便宜。
\end{tipBox}

\textbf{想亲手验证这句话？} \textbf{\texttt{lab09\_flash\_attention/cuda\_softmax/}} 把本节的 Online Softmax 写成真正的 CUDA C++ kernel，在你自己显卡上量出：三个实现\textbf{都贴着 380 GB/s 的 Roofline}（也就是都"优化到位了"），差别只在搬运了 268 MB 还是 134 MB —— 约 1.9x。里面还有一个反直觉的结果：把整行缓存进 shared memory 的版本在 $N = 8192$ 时\textbf{反而掉到 74\%}，因为 shared 占用把 SM 的常驻 block 数从 6 压到 2。\textbf{这正是上面说的"分块"为什么必须存在。}

\vspace{0.8em}\hrule\vspace{0.8em}

\section{[组件] 四、FlashDecoding：把同一个技巧用在 Decode 上}

\subsection{问题}

Decode 时每个请求的 Q 只有 \textbf{1 行}。FlashAttention 靠"不同的 Q 块分给不同的 SM"来并行，现在只剩 batch $\times$ 头数 这么多个独立任务。batch = 1、32 个头、上下文 64K 时，只有 32 个任务，而 H100 有 132 个 SM——大部分 SM 在闲着，每个任务却要顺序扫完 64K 个 key。

\subsection{解法：沿 K/V 的序列维度切分（split-K）}

把长度为 $S$ 的 KV Cache 切成 $s$ 段，\textbf{每段独立}算出自己的 $(m^{(r)}, \ell^{(r)}, \text{acc}^{(r)})$，分给不同的 SM 并行。最后用 Online Softmax 的规则\textbf{合并}：

\[
m = \max_r m^{(r)}, \qquad \ell = \sum_r \ell^{(r)} e^{m^{(r)} - m}, \qquad \text{acc} = \sum_r \text{acc}^{(r)} e^{m^{(r)} - m}, \qquad o = \text{acc} / \ell
\]

\textbf{为什么能这样合并？} 因为第三节的定理说明 $\text{acc}^{(r)} = \sum_{i \in \text{段}r} e^{x_i - m^{(r)}} v_i$，乘以 $e^{m^{(r)} - m}$ 后正好变成 $\sum_{i \in \text{段}r} e^{x_i - m} v_i$，所有段相加就是完整的和。
这个合并操作\textbf{满足结合律}，所以切成几段、按什么顺序合并都不影响结果——这也是为什么它能用在分布式场景（例如把超长上下文的 KV 分到多张卡上，Ring Attention）。

\subsection{与 PagedAttention 的关系}

PagedAttention（第 12 篇）把 KV Cache 分成固定大小的物理块，块在显存里\textbf{不连续}。它的 kernel 本质上就是"一块一块读 KV、用 Online Softmax 累加"——\textbf{Online Softmax 让注意力天然可以分块计算，这是分页 KV Cache 能高效工作的数学前提}。

\vspace{0.8em}\hrule\vspace{0.8em}

\section{五、版本演进一览}

\begin{center}
\small
\begin{tabularx}{\textwidth}{p{3cm} p{4.5cm} X}
\toprule
\textbf{版本} & \textbf{年份} & \textbf{关键改进} \\
\midrule
Online Softmax（Milakov \& Gimelshein） & 2018 & 一遍扫描求 Softmax 的归一化因子 \\
FlashAttention & 2022 & 分块 + Online Softmax + 反向重计算，IO 感知 \\
FlashAttention-2 & 2023 & 交换内外循环、减少非矩阵乘运算、更好的并行划分，约 2 倍提速 \\
FlashDecoding & 2023 & Decode 时沿 KV 序列切分并行 \\
FlashAttention-3 & 2024 & 针对 Hopper：异步数据搬运（TMA）、warp 专门化、FP8 \\
\bottomrule
\end{tabularx}
\end{center}

在 PyTorch 里，\texttt{F.scaled\_dot\_product\_attention} 会自动选择 FlashAttention 等融合 kernel；vLLM、SGLang 则使用 FlashAttention / FlashInfer 等专门的推理注意力库。

\vspace{0.8em}\hrule\vspace{0.8em}

\section{[OK] 本篇自测}

\begin{enumerate}
  \item 朴素注意力的显存和 HBM 访存量各是多少量级？为什么 $d$ 很小时问题特别严重？
  \item 写出 Online Softmax 的更新公式，并用归纳法证明 $\ell_j = \sum_{i \le j} e^{x_i - m_j}$。
  \item 修正因子 $e^{m_{\text{old}} - m_{\text{new}}}$ 为什么永远不会溢出？
  \item FlashAttention 的 FLOPs 比朴素实现少吗？那它为什么快？
  \item FlashDecoding 解决的是什么场景下的什么问题？写出两段部分结果的合并公式。
  \item 为什么说 Online Softmax 是 PagedAttention 的数学前提？
\end{enumerate}

\vspace{0.8em}\hrule\vspace{0.8em}

\section{[资料] 外部资源}

\begin{itemize}
  \item Zihao Ye：\href{https://courses.cs.washington.edu/courses/cse599m/23sp/notes/flashattn.pdf}{From Online Softmax to FlashAttention}（华盛顿大学 CSE 599M 讲义，推导最清晰的短文，\textbf{必读}）
  \item Milakov \& Gimelshein 2018：\href{https://arxiv.org/abs/1805.02867}{Online normalizer calculation for softmax}
  \item Dao et al. 2022：\href{https://arxiv.org/abs/2205.14135}{FlashAttention}；Dao 2023：\href{https://arxiv.org/abs/2307.08691}{FlashAttention-2}；Shah et al. 2024：\href{https://arxiv.org/abs/2407.08608}{FlashAttention-3}
  \item Stanford CRFM 博客：\href{https://crfm.stanford.edu/2023/10/12/flashdecoding.html}{Flash-Decoding for long-context inference}
  \item Triton 官方教程：\href{https://triton-lang.org/main/getting-started/tutorials/06-fused-attention.html}{Fused Attention}（用 Python 写出真正在 GPU 上跑的 FlashAttention；先做完前面的矩阵乘教程）
  \item 代码：\href{https://github.com/Dao-AILab/flash-attention}{Dao-AILab/flash-attention}；\href{https://github.com/gpu-mode/lectures}{GPU MODE 讲座} 中关于 FlashAttention 的几讲
\end{itemize}

\vspace{0.8em}\hrule\vspace{0.8em}

\textbf{上一篇}：\textbf{第 8 篇：硬件视角：GPU 存储层级、FLOPs、算术强度与浮点格式} ｜ \textbf{下一篇}：\textbf{第 10 篇：投机解码：无偏性证明与期望加速比} ｜ \textbf{总路线}：\textbf{LEARNING\_PATH.md}
